% GENERATED FILE. Edit canonical lessons, then run npm run generate:books.
% canonical-source-hash: fnv1a64:3af80597d1ab9268
% canonical-lessons: AR-C343-mughamara, AR-C343-majhul, AR-C343-mufajaa, AR-C343-sudfa, AR-C343-khatara

\chapter{Adventures}
\label{ch:ar-adventures}
\hlchaptermodality{343}

\begin{chapteropening}
\textbf{By the end of this chapter:} \emph{I can say an adventure, unknown, a surprise, a coincidence and to take a risk.}

Complete the last lesson of chapter 343: I can say an adventure, unknown, a surprise, a coincidence and to take a risk.
\end{chapteropening}

\section[\texorpdfstring{\ar{مغامرة} (mughāmara)}{مغامرة (mughāmara)}]{\ar{مغامرة} (mughāmara) — an adventure}

\label{lesson:AR-C343-mughamara}



\begin{quote}
Before the new one: say the Arabic for a title, then the Arabic for to invent.
\end{quote}

\subsection*{You'll want to know: \ar{مغامرة}}
\textbf{\ar{مغامرة}} — \emph{mughāmara} — \textquotedblleft{}an adventure\textquotedblright{}.

\begin{quote}
\ar{كانت} \ar{الرحلة} \ar{مغامرة} \ar{كبيرة}.
\end{quote}

\begin{quote}
\emph{kānat ar-riḥla mughāmara kabīra.}
\end{quote}

\begin{quote}
The trip was a great adventure.
\end{quote}

\begin{cousinweb}
\textbf{\ar{مغامرة}} (\emph{mughāmara}) is built on the native Arabic root \emph{gh-m-r} (Lane), where \emph{ghamr} is deep water that covers a man. \emph{Mughāmara} is the verbal noun of Form III, \textquotedblleft{}to plunge into\textquotedblright{}.
\end{cousinweb}

\begin{grammarlens}[title={What you've built}]
One more word for telling what happened.
\end{grammarlens}

\subsection*{Your turn}
\begin{itemize}
\raggedright
  \item \emph{Say it:} \emph{mughāmara}
  \item \emph{Say it:} \emph{mughāmara}, once more
  \item \emph{Say it:} \emph{kānat ar-riḥla mughāmara kabīra.} — The trip was a great adventure.
  \item \emph{Recall:} say the Arabic for a title, then the Arabic for to invent, then say \emph{mughāmara} again
\end{itemize}

\begin{tcolorbox}[breakable,colback=teal!4,colframe=teal!35!black,title={Before you move on}]
What does \ar{مغامرة} mean? (\textquotedblleft{}An adventure\textquotedblright{}.) Say it once more.
\end{tcolorbox}
\section[\texorpdfstring{\ar{مجهول} (majhūl)}{مجهول (majhūl)}]{\ar{مجهول} (majhūl) — unknown; a stranger}

\label{lesson:AR-C343-majhul}



\begin{quote}
Before the new one: say the Arabic for to invent, then the Arabic for an adventure.
\end{quote}

\subsection*{You'll want to know: \ar{مجهول}}
\textbf{\ar{مجهول}} — \emph{majhūl} — \textquotedblleft{}unknown; a stranger\textquotedblright{}.

\begin{quote}
\ar{وصلت} \ar{رسالة} \ar{من} \ar{شخص} \ar{مجهول}.
\end{quote}

\begin{quote}
\emph{waṣalat risāla min shakhṣ majhūl.}
\end{quote}

\begin{quote}
A letter arrived from an unknown person.
\end{quote}

In grammar, the passive is \ar{المبني} \ar{للمجهول}, built for the unknown doer.

\begin{cousinweb}
\textbf{\ar{مجهول}} (\emph{majhūl}) is usually traced to the native Arabic root \emph{j-h-l}, \textquotedblleft{}to be ignorant, not to know\textquotedblright{}. \emph{Majhūl} is the passive participle: \textquotedblleft{}what is not known\textquotedblright{}.
\end{cousinweb}

\begin{grammarlens}[title={What you've built}]
One more word for telling what happened.
\end{grammarlens}

\subsection*{Your turn}
\begin{itemize}
\raggedright
  \item \emph{Say it:} \emph{majhūl}
  \item \emph{Say it:} \emph{majhūl}, once more
  \item \emph{Say it:} \emph{waṣalat risāla min shakhṣ majhūl.} — A letter arrived from an unknown person.
  \item \emph{Recall:} say the Arabic for to invent, then the Arabic for an adventure, then say \emph{majhūl} again
\end{itemize}

\begin{tcolorbox}[breakable,colback=teal!4,colframe=teal!35!black,title={Before you move on}]
What does \ar{مجهول} mean? (\textquotedblleft{}Unknown; a stranger\textquotedblright{}.) Say it once more.
\end{tcolorbox}
\section[\texorpdfstring{\ar{مفاجأة} (mufājaʾa)}{مفاجأة (mufājaʾa)}]{\ar{مفاجأة} (mufājaʾa) — a surprise}

\label{lesson:AR-C343-mufajaa}



\begin{quote}
Before the new one: say the Arabic for an adventure, then the Arabic for unknown.
\end{quote}

\subsection*{You'll want to know: \ar{مفاجأة}}
\textbf{\ar{مفاجأة}} — \emph{mufājaʾa} — \textquotedblleft{}a surprise\textquotedblright{}.

\begin{quote}
\ar{كانت} \ar{الهدية} \ar{مفاجأة} \ar{جميلة}.
\end{quote}

\begin{quote}
\emph{kānat al-hadiyya mufājaʾa jamīla.}
\end{quote}

\begin{quote}
The present was a lovely surprise.
\end{quote}

\begin{cousinweb}
\textbf{\ar{مفاجأة}} (\emph{mufājaʾa}) is built on the native Arabic root \emph{f-j-ʾ} (Lane), \textquotedblleft{}to come upon unawares\textquotedblright{}, the root of \emph{fajʾatan}, suddenly. It is the verbal noun of Form III.
\end{cousinweb}

\begin{grammarlens}[title={What you've built}]
One more word for telling what happened.
\end{grammarlens}

\subsection*{Your turn}
\begin{itemize}
\raggedright
  \item \emph{Say it:} \emph{mufājaʾa}
  \item \emph{Say it:} \emph{mufājaʾa}, once more
  \item \emph{Say it:} \emph{kānat al-hadiyya mufājaʾa jamīla.} — The present was a lovely surprise.
  \item \emph{Recall:} say the Arabic for an adventure, then the Arabic for unknown, then say \emph{mufājaʾa} again
\end{itemize}

\begin{tcolorbox}[breakable,colback=teal!4,colframe=teal!35!black,title={Before you move on}]
What does \ar{مفاجأة} mean? (\textquotedblleft{}A surprise\textquotedblright{}.) Say it once more.
\end{tcolorbox}
\section[\texorpdfstring{\ar{صدفة} (ṣudfa)}{صدفة (ṣudfa)}]{\ar{صدفة} (ṣudfa) — a coincidence, chance}

\label{lesson:AR-C343-sudfa}



\begin{quote}
Before the new one: say the Arabic for unknown, then the Arabic for a surprise.
\end{quote}

\subsection*{You'll want to know: \ar{صدفة}}
\textbf{\ar{صدفة}} — \emph{ṣudfa} — \textquotedblleft{}a coincidence, chance\textquotedblright{}.

\begin{quote}
\ar{التقيت} \ar{بصديقي} \ar{القديم} \ar{صدفة} \ar{في} \ar{السوق}.
\end{quote}

\begin{quote}
\emph{iltaqaytu bi-ṣadīqī al-qadīm ṣudfatan fī as-sūq.}
\end{quote}

\begin{quote}
I met my old friend by chance in the market.
\end{quote}

\ar{صدفة} or \ar{بالصدفة}: by chance.

\begin{cousinweb}
\textbf{\ar{صدفة}} (\emph{ṣudfa}) is built on the native Arabic root \emph{ṣ-d-f} (Lane), where Form III \emph{ṣādafa} is \textquotedblleft{}to light upon someone without meaning to\textquotedblright{}.
\end{cousinweb}

\begin{grammarlens}[title={What you've built}]
One more word for telling what happened.
\end{grammarlens}

\subsection*{Your turn}
\begin{itemize}
\raggedright
  \item \emph{Say it:} \emph{ṣudfa}
  \item \emph{Say it:} \emph{ṣudfa}, once more
  \item \emph{Say it:} \emph{iltaqaytu bi-ṣadīqī al-qadīm ṣudfatan fī as-sūq.} — I met my old friend by chance in the market.
  \item \emph{Recall:} say the Arabic for unknown, then the Arabic for a surprise, then say \emph{ṣudfa} again
\end{itemize}

\begin{tcolorbox}[breakable,colback=teal!4,colframe=teal!35!black,title={Before you move on}]
What does \ar{صدفة} mean? (\textquotedblleft{}A coincidence, chance\textquotedblright{}.) Say it once more.
\end{tcolorbox}
\section[\texorpdfstring{\ar{خاطر} (khāṭara)}{خاطر (khāṭara)}]{\ar{خاطر} (khāṭara) — to take a risk}

\label{lesson:AR-C343-khatara}



\begin{quote}
Before the new one: say the Arabic for a surprise, then the Arabic for a coincidence.
\end{quote}

\subsection*{You'll want to know: \ar{خاطر}}
\textbf{\ar{خاطر}} — \emph{khāṭara} — \textquotedblleft{}to take a risk\textquotedblright{}.

\begin{quote}
\ar{خاطر} \ar{الرجل} \ar{بحياته} \ar{لينقذ} \ar{الطفل}.
\end{quote}

\begin{quote}
\emph{khāṭara ar-rajul bi-ḥayātihi li-yunqidha aṭ-ṭifl.}
\end{quote}

\begin{quote}
The man risked his life to save the child.
\end{quote}

Present: \ar{يخاطر}, \emph{yukhāṭir}.

It takes \ar{بـ} for what is risked.

\begin{cousinweb}
\textbf{\ar{خاطر}} (\emph{khāṭara}) is built on the native Arabic root \emph{kh-ṭ-r} (Lane), the root of \emph{khaṭar}, danger, and \emph{khaṭīr}, serious. It is Form III, \textquotedblleft{}to put a thing in danger\textquotedblright{}.
\end{cousinweb}

\begin{grammarlens}[title={What you've built}]
One more word for telling what happened.
\end{grammarlens}

\subsection*{Your turn}
\begin{itemize}
\raggedright
  \item \emph{Say it:} \emph{khāṭara}
  \item \emph{Say it:} \emph{khāṭara}, once more
  \item \emph{Say it:} \emph{khāṭara ar-rajul bi-ḥayātihi li-yunqidha aṭ-ṭifl.} — The man risked his life to save the child.
  \item \emph{Recall:} say the Arabic for a surprise, then the Arabic for a coincidence, then say \emph{khāṭara} again
\end{itemize}

\begin{tcolorbox}[breakable,colback=teal!4,colframe=teal!35!black,title={Before you move on}]
What does \ar{خاطر} mean? (\textquotedblleft{}To take a risk\textquotedblright{}.) Say it once more.
\end{tcolorbox}
